\begin{frame}{What the badge model cannot establish}
\princ{MODEL LIMIT}
The two-state model can establish only claims about its declared states,
inputs, transitions, start state, and trace.
\begin{itemize}
\item It does not establish whether a person was authorized.
\item It does not establish whether a physical reader worked.
\item It does not represent denied, expired, or tampered conditions.
\end{itemize}
\begin{block}{Strong conclusion}
The trace is correct for the declared model. It is not a proof that a real access system is safe.
\end{block}
\see{a limitation is a precise omitted condition, not a vague disclaimer.}
\end{frame}

\begin{frame}{The three AI failure modes}
\princ{FAILURE MODES}
\begin{tabular}{@{}p{4.3cm}p{4.4cm}p{3.6cm}@{}}
\textbf{Failure} & \textbf{Check} & \textbf{Visible evidence}\\\midrule
False simplification & all truth-table rows & one counterexample row\\[4pt]
Witness called universal proof & separate SAT/tautology tests & witness plus remaining counterexample\\[4pt]
Omitted state or transition & declared Cartesian product & named missing pair
\end{tabular}
\medskip
\predict{Which check would you demand first when an answer sounds too certain?}
\see{choose the check that matches the claim: rows, witnesses, or transitions.}
\end{frame}

\begin{frame}{Tuesday individual activity}
\princ{MAKE YOUR OWN EVIDENCE}
\begin{enumerate}
\item Complete the supplied $E$ versus $S$ truth table.
\item Give one satisfying assignment and explain why it is not a universal proof.
\item Trace valid, timeout, valid from Locked.
\item Name one condition the two-state model cannot establish.
\item Write one sentence separating what the evidence proves from what it leaves open.
\end{enumerate}
\see{work individually first; if AI is used, record the exact row or transition you checked.}
\end{frame}

\begin{frame}{Preview: defend a bounded model}
\princ{EVIDENCE BEFORE CONFIDENCE}
On Thursday, choose one path:
\begin{itemize}
\item Boolean/circuit: expression, circuit or simplification, and complete table; or
\item state machine: states, inputs, transitions, start state, two traces, and a limit.
\end{itemize}
\begin{block}{Defense norm}
State the claim, show the row or trace, name the strongest limitation, answer a challenge, and critique the evidence rather than the person.
\end{block}
\see{a narrowed claim is a successful result when the challenge earns it.}
\end{frame}

\begin{frame}{Exit sentence}
\princ{LIMIT THE CLAIM}
\begin{block}{Complete this sentence}
For this declared model, the evidence establishes \underline{\hspace{4cm}}; it does not establish \underline{\hspace{4cm}} because \underline{\hspace{3cm}}.
\end{block}
\begin{itemize}
\item A truth table establishes behavior for the assignments it exhausts.
\item A trace establishes the path for the inputs and transitions it names.
\item A model limit protects the claim from becoming larger than its evidence.
\end{itemize}
\note{Point students toward the canonical Checkpoint 3 source for Thursday; this deck does not reproduce that source.}
\end{frame}
